\documentclass[11pt]{article}
\usepackage[a4paper,margin=1in]{geometry}
\usepackage{amsmath,amssymb}
\usepackage{booktabs}
\usepackage{microtype}
\usepackage{xcolor}
\usepackage{titlesec}
\usepackage{enumitem}
\usepackage[hidelinks]{hyperref}

\definecolor{accent}{RGB}{15,110,86}
\definecolor{amber}{RGB}{186,117,23}
\definecolor{coral}{RGB}{153,60,29}

\titleformat{\section}{\color{accent}\large\bfseries}{\thesection}{0.6em}{}
\titleformat{\subsection}{\color{accent!85!black}\normalsize\bfseries}{\thesubsection}{0.6em}{}

\newcommand{\dcor}{\mathrm{dCor}}
\newcommand{\pdcor}{\mathrm{pdCor}}
\newcommand{\I}{I}
\newcommand{\KL}{D_{\mathrm{KL}}}
\newcommand{\keybox}[1]{%
  \begin{center}\fcolorbox{accent}{accent!6}{%
  \begin{minipage}{0.93\linewidth}\vspace{4pt}#1\vspace{4pt}\end{minipage}}\end{center}}

\title{\vspace{-2.2em}\color{accent}\bfseries
A Two-Object Dependence Pipeline for the LLP ABCD Analysis}
\author{\normalsize Nathan Herling \\[-2pt]
\small UA ATLAS HEP Group \,$\cdot$\, Project HELIX \,$\cdot$\, Run~3 LLP MSVtx search}
\date{\small\today}

\begin{document}
\maketitle
\vspace{-1.2em}\hrule height 1.2pt \vspace{1em}

\noindent
Prof.\ Johns's question has two distinct objects hiding inside it. ``Are the
features the nets consist of uncorrelated?'' is about \emph{inputs}; ``do the
two discriminants behave independently for ABCD?'' is about \emph{outputs}.
They are linked but not identical, and answering one does not answer the other.
This note lays out a pipeline that measures both with distribution-free metrics,
connects them through a single theorem, and assigns a concrete interpretation to
each result pattern. A companion note argues why Pearson is the wrong metric and
distance correlation (dCor) the right one; that argument is assumed here.

\keybox{\centering \textbf{The spine:} inputs set a \emph{ceiling} on how
dependent the outputs can be; the outputs are the \emph{realized} value;
conditioning tells you whether output dependence was \emph{inherited} or
\emph{manufactured}; the slice tells you \emph{where} it lives.}

\section{The theorem that links the two objects}
Both nets are deterministic maps of their inputs, $s_1=f(X_1)$, $s_2=g(X_2)$.
The data-processing inequality gives a hard ceiling,
\begin{equation}
\I(s_1;s_2)\;\le\;\I(X_1;X_2).
\end{equation}
The networks cannot manufacture more mutual information between their outputs
than already exists between their inputs. Two consequences:
\begin{itemize}[nosep]
\item independent input sets ($\I(X_1;X_2)=0$) \emph{force} independent outputs
      --- this is exactly why ``ideally the features are uncorrelated'' is the
      right instinct;
\item the bound is loose: dependent inputs only \emph{permit} dependent outputs.
      The nets may discard the shared information and sit far below the ceiling.
\end{itemize}
So input dependence is necessary-but-not-sufficient for output dependence, and
the \emph{gap} between ceiling and realized value is the interesting physics.
(The clean inequality is for mutual information; dCor does not obey an equally
tight processing inequality, so MI carries the \emph{theorem} and dCor/HSIC
carry the \emph{measurement} --- they agree on the iff-zero, which is what the
independence question needs.)

\section{Object A --- input feature correlation}
\subsection{A1 : within-set redundancy}
For each net's own inputs: the Pearson and dCor correlation matrices (linear vs.\
nonlinear redundancy), the PCA eigenvalue spectrum summarised by the
\emph{participation ratio} $\mathrm{PR}=(\sum_i\lambda_i)^2/\sum_i\lambda_i^2$
(effective input dimensionality), and the Gaussian \emph{total correlation}
$C(X)=-\tfrac12\log\det R$ as a single redundancy number. A low PR means fewer
effective inputs than nominal --- the set is prunable and the net is more prone
to converging on a shared representation.

\subsection{A2 : cross-set overlap (the ceiling)}
This is the object Johns named. It needs \emph{set-to-set} measures, not pairwise:
\begin{itemize}[nosep]
\item Jaccard overlap of the literal feature lists --- shared raw variables are
      the most obvious dependence channel;
\item canonical correlation analysis (CCA) --- the classical linear spectrum,
      its leading value the max linear correlation between linear combinations
      of the two sets;
\item \textbf{vector dCor} $\dcor(X_1\text{-vec},X_2\text{-vec})\in[0,1]$ ---
      dCor is defined between random \emph{vectors}, so a single number captures
      the nonlinear set-to-set dependence; this is the headline A2 metric;
\item normalized HSIC --- the kernel cousin of vector dCor, with its own
      permutation test, run as an estimator cross-check.
\end{itemize}

\section{Object B --- output score correlation}
The ladder from the companion note, now as a pipeline: Pearson (a floor only) and
Spearman; dCor with a permutation test; mutual information $\I(s_1;s_2)$ (the
realized value to compare against the A2 ceiling); sliced dCor across NN1 plus
the conditional mean \emph{and variance} profiles $\mathbb{E}[s_2\mid s_1]$,
$\mathrm{Var}(s_2\mid s_1)$ (locality and the heteroscedastic channel); and the
bootstrapped ABCD closure ratio as the final physics arbiter.

\section{The bridge --- inherited vs.\ manufactured}
The piece that turns two measurements into an explanation is the
\emph{partial distance correlation} (Sz\'ekely--Rizzo), $\pdcor(s_1,s_2;Z)$,
which projects the shared features $Z$ out of the score--score dependence:
\begin{equation}
\pdcor(s_1,s_2;Z)\approx 0 \;\Rightarrow\; \text{dependence \textbf{inherited}
(feature-level fix viable)},
\end{equation}
\begin{equation}
\pdcor(s_1,s_2;Z)\ \text{survives} \;\Rightarrow\; \text{dependence
\textbf{manufactured} (output-level fix or MaxEnt tilt)}.
\end{equation}
Its rigorous complement, requiring the trained models, is \emph{feature-level
event mixing}: rebuild the two nets' inputs from different events to destroy the
shared latent and re-measure output dCor; the drop quantifies the share of
output dependence due to the common event. This is the same operation as the
event-mixing background estimate, repurposed as a diagnostic.

\section{Result interpretations, from theory alone}
\begin{center}
\renewcommand{\arraystretch}{1.3}
\begin{tabular}{@{}p{4.6cm}p{9.8cm}@{}}
\toprule
\textbf{Pattern} & \textbf{Reading} \\
\midrule
high input cross-dCor + high output dCor &
inherited baseline; outputs near the ceiling. ABCD still suffers, but the fix is
upstream (more orthogonal inputs). \\
low input cross-dCor + high output dCor &
multivariate shared latent the pairwise scan missed; feature decorrelation
futile, go to output-level (DisCo) or model it (MaxEnt tilt). \\
high input cross-dCor + low output dCor &
the good case: nets learned to decorrelate. Verify closure, then relax. \\
Pearson $\approx$ dCor (outputs) &
dependence mostly linear $\Rightarrow$ near-Gaussian joint $\Rightarrow$ a single
MaxEnt tilt likely captures the bias; barrel may be cheaply fixable. \\
Pearson $\ll$ dCor (e.g.\ $0.3$ vs $0.6$) &
strongly nonlinear; $0.3$ was a bad floor. Bias structured, likely tail-local;
a single tilt may not suffice. \\
sliced dCor flat vs.\ spiking in top corner &
flat $\Rightarrow$ global correction works; spike at the signal corner
$\Rightarrow$ local closure failure even at modest global dCor (prior for the
barrel $1.79$). \\
$\mathrm{Var}(s_2\mid s_1)$ varies, $\mathbb{E}[s_2\mid s_1]$ flat &
pure variance dependence: $\rho\approx0$ yet ABCD breaks; mean-based corrections
fail. Only dCor / conditional-variance scan catches it. \\
low within-set PR &
fewer effective inputs than nominal; predicts high output dependence and flags
the feature set as prunable. \\
\bottomrule
\end{tabular}
\end{center}

\keybox{This four-part story --- ceiling, realized value, inherited-vs-manufactured,
and locality --- is a stronger answer than any single correlation number: it
explains not just \emph{whether} the discriminants are dependent, but \emph{why}
and \emph{where}.}

\section{Software}
The accompanying suite (\texttt{run\_diagnostics.py}) implements all of the
above: \texttt{core\_metrics.py} (dCor, permutation, HSIC, MI, partial dCor,
CCA), \texttt{feature\_module.py} (Object A), \texttt{output\_module.py}
(Object B), \texttt{bridge\_module.py} (partial dCor + event-mixing hook). Run
\texttt{python run\_diagnostics.py -{}-demo} for a synthetic shared-latent
example, or point it at the ROOT$\to$CSV exports via \texttt{config.py}. All
dependence metrics are evaluated on subsamples (dCor/HSIC are $O(n^2)$ in
memory); run on background-only, continuous pre-quantization scores.

\vspace{0.6em}\hrule height 0.5pt \vspace{0.4em}
{\footnotesize\itshape The pipeline establishes which channel dominates ---
input overlap, manufactured coupling, tail localisation, or variance dependence
--- but the constancy of any MaxEnt tilt across the plane (control region to
signal region) remains an empirical claim to validate, not prove. That
extrapolation is where most of the residual risk lives.}

\end{document}
